\subsection{Sprint 0}

Na Sprint 0, concentrei minha atuação na estruturação inicial do projeto. Comecei pela análise do termo de abertura, buscando
compreender o problema proposto, o escopo inicial e a viabilidade de realizá-lo
dentro do cronograma do semestre. A partir dessa leitura, formulei possíveis
perguntas para o levantamento de requisitos, que seriam discutidas posteriormente
com o restante do time.

Também participei da organização do ambiente de trabalho e da equipe. Criei e
realizei ajustes no servidor de Discord utilizado para a comunicação do projeto.
Em uma discussão sobre a forma de organização do desenvolvimento, avaliamos a
divisão da equipe por camada técnica, entre \textit{frontend} e \textit{backend}, e
a formação de \textit{squads}. Considerando a experiência, o engajamento e a
distribuição das responsabilidades entre os integrantes, optamos inicialmente por
trabalhar em \textit{squads}, deixando aberta a possibilidade de rever essa decisão
após a primeira sprint. Essa organização exigiu um alinhamento maior do que o
habitual, pois o projeto contava com cinco alunos AGES IV,
incluindo a mim.

Após a reunião inicial com as \textit{stakeholders}, surgiram divergências sobre a
ordem das atividades de definição do produto. Enquanto parte do grupo entendia
que seria necessário fechar o escopo antes de avançar, outra parte defendia o
início dos esboços das telas. A falta de consenso atrasou consideravelmente a
elaboração dos protótipos e tornou as referências iniciais menos claras para os
integrantes que ainda estavam se familiarizando com o projeto.

Na etapa de preparação da primeira entrega, participei da finalização dos
protótipos das telas e da preparação da apresentação para as
\textit{stakeholders}. A apresentação abordou as telas e as histórias de usuário
planejadas para a sprint seguinte. Após esse alinhamento, também elaborei e
organizei as tarefas no quadro do projeto para a Sprint 1, transformando o que
havia sido discutido com as \textit{stakeholders} em itens de trabalho para o
time.
\\

\fbox{%
    \parbox{0.95\linewidth}{%
        \setlength{\parindent}{15pt}
        \itshape A composição do time me deixou receoso desde o início. Em minhas
        experiências anteriores, era mais comum haver poucos AGES III e AGES IV,
        muitas vezes trabalhando em duplas. Esse formato facilitava tanto a
        referência técnica para \textit{frontend}, \textit{backend} e infraestrutura
        quanto as decisões de gestão entre os responsáveis pelo projeto. Nesta
        experiência, porém, havia cinco alunos em cada uma dessas etapas, o que
        tornou difícil chegar a um acordo mesmo sobre decisões iniciais.

        O ponto que mais me desgastou na Sprint 0 foi perceber que, após a reunião
        com as \textit{stakeholders}, o time permaneceu dividido entre fechar o
        escopo e começar os esboços das telas. O atraso dos protótipos foi uma
        consequência direta desse impasse e, no meio dele, tive a impressão de que
        os colegas iniciantes ficavam ainda mais perdidos, justamente quando
        precisavam de direcionamento.

        Essa situação me atingiu de forma particular porque, na AGES III, já havia
        passado por conflitos no início de um projeto que reduziram bastante meu
        interesse por ele. Ao encontrar novamente um ambiente de discordância, quase
        perdi a paciência com o projeto como um todo. A Sprint 0 reforçou para mim
        que ter mais pessoas em papéis de liderança não torna a organização mais
        fácil por si só: sem uma forma clara de decidir e seguir adiante, a discussão
        pode paralisar o trabalho e prejudicar principalmente quem ainda está
        aprendendo a atuar no time.
    }%
}
